% Chute libre : énergies cinétique, potentielle et mécanique en trois points (L'essentiel, chapitre 16)
% Balle de 0,60 kg lâchée de 2,00 m : Em = 11,8 J
\begin{tikzpicture}[x=1cm,y=1cm]
% Axe et positions
\draw[otInk,line width=0.6pt,-{Stealth[length=2mm]}] (0,-0.2) -- (0,3.3) node[pcLegende,anchor=south] {$z$};
\fill[otMuted!25] (-0.3,-0.25) rectangle (1.6,0);
\draw[otInk!60,line width=0.5pt] (-0.3,0) -- (1.6,0);
\node[pcLegende,anchor=east] at (-0.05,0) {$z=0$};
\foreach \z/\n in {3.0/A,1.8/B,0.35/C}{%
  \shade[ball color=otOrange!60] (0.8,\z) circle (0.13);
  \node[pcLegendeG,anchor=west] at (0.98,\z) {\n};}
\draw[otMuted,line width=0.4pt,dash pattern=on 1.5pt off 1.2pt] (0.8,2.85) -- (0.8,0.5);
% Diagrammes en barres
\foreach \bx/\ep/\eco/\n in {3.0/2.8/0/A,5.6/1.68/1.12/B,8.2/0.33/2.47/C}{%
  \begin{scope}[shift={(\bx,0)}]
    \fill[otNavy!55] (0,0) rectangle (0.5,\ep);
    \fill[otOrange!70] (0.65,0) rectangle (1.15,\eco);
    \fill[otGreenDark!70] (1.3,0) rectangle (1.8,2.8);
    \draw[otInk!60,line width=0.5pt] (-0.15,0) -- (1.95,0);
    \node[pcLegende,anchor=north] at (0.25,-0.02) {$E_\mathrm{p}$};
    \node[pcLegende,anchor=north] at (0.9,-0.02) {$E_\mathrm{c}$};
    \node[pcLegende,anchor=north] at (1.55,-0.02) {$E_\mathrm{m}$};
    \node[pcLegendeG,anchor=south] at (0.9,3.0) {point \n};
  \end{scope}}
\end{tikzpicture}
\par\smallskip\legende{Sans frottement, l'énergie potentielle perdue se retrouve en énergie cinétique : l'énergie mécanique reste constante.}
